\documentclass[a4paper]{article} %

\date{} %

% Please, do not change any of the above lines

\usepackage{amssymb} % add other LaTeX packages you need here.

\begin{document}
\setcounter{page}{1} % Please, do not delete this line

% Your title
\title{On representations of double Boolean algebras}

% Authors
\author{Leonard Kwuida \\ %
       Bern University of Applied Sciences\\ \\ % Affiliation 1
      % Author 4 \\ % If any other author with different Affilation
       %University of ZZWW\\ \\ % Affiliation 2 (if needed)
       % New author \\
       % New affiliation \\
       % Add authors and affiliation as needed
       \texttt{leonard.kwuida©bfh.ch} % Only one corresponding e-mail, please
       }%

\maketitle

% The abstract

Double Boolean algebras are abstract structures that arise in Formal Concept Analysis from the need to develop a Boolean Concept Logic, based on concepts as units of thought. Several representations have been addressed including contextual (pre-, semi-, proto-concept algebras)\cite{HLSW00,Wi00}, topological and logical\cite{HB23,HKBL26,TKKTT2025} representations. We will review some of these in our presentation. This is an effort of the presenter to understand what has been done so far, avoiding reinventing the wheel.  

\begin{thebibliography}{99}
\bibitem{HLSW00}        
C. Herrmann, P. Luksch, M. Skorsky, R. Wille. \newblock \emph{Algebras of semiconcepts and double Boolean algebras}. In: Contributions to General Algebra \textbf{13}. Verlag Johannes heyn, Klagenfurt 2001, 175-188

\bibitem{HB23}
P. Howlader, M. Banerjee. \newblock \emph{Topological representation of double Boolean algebras}. Algebra Univers. \textbf{84}~15, 2023 doi:10.1007/s00012-023-00811-x

\bibitem{HKBL26} 
P. Howlader, L. Kwuida, M. Behrisch, C-J. Liau. \newblock \emph{Towards a Simplified Theory of Double Boolean Algebras: Axioms and Topological Representation}. Preprint doi:10.48550/arXiv.2601.01418
\bibitem{TKKTT2025}
G. Tenkeu K.,  M. Krebs, M., L. Kwuida, E.R. Temgoua A., Y.L. Tenkeu J. \newblock \emph{A Stone Duality for Pure and Trivial Double Boolean Algebras}. 
In: Cellier, P., Ganter, B., Missaoui, R. (eds) Conceptual Knowledge Structures. CONCEPTS 2025. LNCS \textbf{15941}. Springer. doi:10.1007/978-3-032-03364-2\_26
\bibitem{Wi00}
R. Wille. \newblock \emph{Boolean Concept Logic}. In: Ganter, B., Mineau, G.W. (eds) Conceptual Structures: Logical, Linguistic, and Computational Issues. ICCS 2000. LNCS \textbf{1867}. Springer. doi:10.1007/10722280\_22.
      



\end{thebibliography}

\end{document}